% !TEX root = chapter-standalone.tex

\section{Group actions}

\todotext{This to be made its own section, and rewritten in the same way monoid actions are to be rewritten}

For defining a group action, we must adapt the definition of a monoid action.
The endomorphisms~$\Endof\setA$ are not a group, because they also contain non-invertible maps.
Recall that an invertible \SY{endomorphism} is called an \emph{automorphism}, and that~$\Autof\setA$, the set of \SY{automorphisms} of~\setA, comes naturally equipped with a \SY{group} structure.

\todotextjira{730}{\bernina: @JL: better would be to make examples in the monoid and group sections for End(A) and Aut(A), respectively, and then reference those here.}

\begin{ctdefinition}
    [Group right action]\label{def:group-right-action-operation}
    A \emph{right action} of a \SY{group}~$\grpA = \tup{\grpAset, \mtimesof\grpA, \idmon_\grpA}$ on a set~\setA is:

    \constit
    \begin{enumerate}
        \item a function
              \begin{equation}
                  \label{eq:ract-group}
                  \act \colon \makecartprod{\setA, \grpAset} \sto \setA.
              \end{equation}
    \end{enumerate}
    \condit
    \begin{enumerate}
        \item for all~$\setAel \setin \setA$ and $\monela,\monelb\setin\grpAset$,
              \begin{equation}
                  \label{eq:group-ract-1cond}
                  \act(\act(\setAel,\monela), \monelb) = \act(\setAel, \monela \mtimesof\monoidA \monelb);
              \end{equation}
        \item for all~$\setAel \setin \setA$,
              \begin{equation}
                  \label{eq:group-ract-2cond}
                  \act(\setAel, \idmon_\monoidA) = \setAel.
              \end{equation}
    \end{enumerate}
\end{ctdefinition}

\todotext{JL: make remark or lemma here that inverses being respected comes for free... and spell out the property in action terms...}

\begin{ctdefinition}[Group action]
    \label{def:group-cov-action}
    An \maindef{action} of a \SY{group}~$\grpA$ onto a set~\setA is a \SY{group morphism}
    \begin{equation}
        \label{eq:act-group}
        \act \colon \grpA \fto \Autof\setA.
    \end{equation}
\end{ctdefinition}

\todotext{\bernina: @JL: add example of vector spaces as actions of a field on an abelian group; mention concept of a module}

\todotext{lemma commented out here}
%\begin{lemma}
%\label{lem:uncurrying-group-actions}
%The data of a \SY{group action}
%\begin{equation}
%\act \colon \grpA \fto \Autof\setA
%\end{equation}
%is equivalent to the data of having a function
%\begin{equation}
%\stylemaps{\alpha} \colon \grpAset \cartprod \setA \sto \setA,
%\end{equation}
%which satisfies the conditions
%\begin{enumerate}
%\item $ \stylemaps{\alpha} (\idgrp_\grpA, \setAel) = \setAel \quad \forall \setAel \setin \setA$;
%\item $ \stylemaps{\alpha} (\grpelb, \stylemaps{\alpha}(\grpela, \setAel)) = \stylemaps{\alpha} (\grpela \mtimes \grpelb, \setAel) \quad \forall \setAel \setin \setA, \forall \grpela, \grpelb \setin \grpAset.
%      $
%\end{enumerate}
%\end{lemma}


\todotext{Add examples, in particular ones relating to group representations and even some definitions perhaps of matrix groups, understood via actions e.g. when notion of symmetry is in the foreground.}



\begin{example}[Time dilation of signals]
In the field of signal-processing, a basic way to model a signal is as a function $f \colon \reals \to \subA$, where the domain models time and the codomain models the possible values that a signal might take. 

In this example, we consider the set of all possible such signals, $\setA = \subA^\reals$. Let $\tup{\posReals, \cdot, 1}$ be the group of positive reals, with multiplication as composition. We can define a group action $\act \colon \posReals \to \Autof\setA$ on signals that models ``time dilation'':
\begin{equation}
\defmapperiod{\act(\lambda)}{\setA}{\to}{\setA}{f(t)}{f(t / \lambda)}
\end{equation}
\end{example}




\devel{
\subsection{Linear group actions}
\label{sec:action-matrix-groups}

\todotext{J: this section should be moved into the section on group actions }

\paragraph{Special Euclidean group}

\begin{definition}
    [General Euclidean group~$\mgEn$]
    \label{def:general-euclidean-group}
    The general (real) Euclidean group of order~$n$, written~$\mgEn$, is the group of~$\matdim{(n+1)}{(n+1)}$ real matrices of the form
    \begin{equation}
        \label{eq:En-def}
        \Ematrix{\matmor{R}}{\vectmor{t}},
    \end{equation}
    where~$\matmor{R}\setin \mgOn$ and~$\vectmor{t} \setin \reals^n$.
\end{definition}

\begin{definition}
    [Special Euclidean group~$\mgSEn$]
    \label{def:special-euclidean-group}
    The special (real) Euclidean group of order~$n$, written~$\mgSEn$, is the group of~$(n+1)\times (n+1)$ real matrices of the form
    \begin{equation}
        \label{eq:SEn-def}
        \Ematrix{\matmor{R}}{\vectmor{t}},
    \end{equation}
    where~$\matmor{R}\setin \mgSOn$ and~$\vectmor{t} \setin \reals^n$.
\end{definition}

\begin{table*}
    \caption{Matrix groups}
    \label{tab:matrix-groups}
    \begin{tabular}{cllc}
        \mgGLn & General linear group     & invertible linear transformations \\
        \mgSLn & Special linear group     & invertible linear transformations with determinant equal to 1 \\ \\
        \mgOn  & Orthogonal group         & preserve length of vectors \\
        \mgSOn & Special orthogonal group & rotations \\ \\
        \mgEn  & Euclidean group          & preserve distances and angles                                 & \\
        \mgSEn & Special Euclidean group  & rigid motions \\
    \end{tabular}
\end{table*}

The \SY{groups}~$\mgSEtwo$ and~$\mgSEthree$ are particularly important in robotics because they represent the roto-translations of the plane and 3D space, respectively.

From~\cref{eq:SEn-def} we know we can represent one by a pair~$\tupp{\matmor{R}, \vectmor{t}}$, with~$\matmor{R}\setin\mgSOn$ and~$\vectmor{t} \setin \reals^n$.

If we look at how matrices compose, we get
%
\begin{equation}
    \label{eq:SE-composition}
    \Ematrix{\matmor{R}_2}{\vectmor{t}_2} \Ematrix{\matmor{R}_1}{\vectmor{t}_1} = \Ematrix{\matmor{R}_2 \matmor{R}_1}{\matmor{R}_2\vectmor{t}_1 + \vectmor{t}_2}.
\end{equation}
%
The formula for composition is
%
\begin{equation}
    \label{eq:SE-composition-short}
    \tupp{\matmor{R}_1, \vectmor{t}_1} \mtimesof{\mgSEn} \tupp{\matmor{R}_2, \vectmor{t}_2} = \tupp{\matmor{R}_2 \matmor{R}_1, \matmor{R}_2 \vectmor{t}_1 + \vectmor{t}_2}.
\end{equation}
%

% %
% \begin{equation}
%   \tupp{\matmor{R}_1, \vectmor{t}_1} \mthenof{\mgSEn} \tupp{\matmor{R}_2, \vectmor{t}_2}  = \tupp{\matmor{R}_2 \matmor{R}_1, \matmor{R}_2 \vectmor{t}_1 + \vectmor{t}_2}
% \end{equation}
%
% Note that because we are using $\mthen$ (then), we do not have the nice symmetry
% that we had if we were using $\after$ (after).
% %
% %

The group~$\mgSEn$ induces a transformation on the points of~$\reals^n$.
We are going to call this an \emph{action}.

The action is the following function:
%
\begin{equation}
    \label{eq:apply}
    \defmapperiodset{\mgact}{
        \mgSEn \cartprod \reals^n
    }{
        \reals^n
    }{
        \tupp{ \tupp{\matmor{R}, \vectmor{t}}, \vectob{p}}
    }{
        \matmor{R}\vectob{p} + \vectmor{t}
    }
\end{equation}
Given a roto-translation and a point, the function returns the roto-translated point.
%
We can also see this in matrix form as follows.
We need to substitute for a point~$\vectob{p} \setin \reals^n$ a homogenous point~$\Epoint{\vectob{p}} \setin \reals^{n+1}$.
%
\begin{equation}
    \label{eq:homogenous}
    \Ematrix{\matmor{R}}{\vectmor{t}}
    \Epoint{\vectob{p}}
    =
    \Epoint{\matmor{R}\vectob{p} + \vectmor{t}}.
\end{equation}
%

If we apply two rototranslations, first~$\monela=\tupp{\matmor{R}_{\monela}, \vectmor{t}_{\monela}}$ and then~$\monelb=\tupp{\matmor{R}_{\monelb}, \vectmor{t}_{\monelb}}$, we find:
%
\begin{equation}
    \label{eq:rototranslation-series}
    \begin{aligned}
        \mgact(\tupp{\matmor{R}_{\monelb}, \vectmor{t}_{\monelb}}, \mgact(\tupp{\matmor{R}_{\monela}, \vectmor{t}_{\monela}}, \vectob{p})) & =
        \mgact(\tupp{\matmor{R}_{\monelb}, \vectmor{t}_{\monelb}}, \matmor{R}_{\monela}\vectob{p} + \vectmor{t}_{\monela}) \\
                                                                                                                                           & = \matmor{R}_{\monelb}\matmor{R}_{\monela}\vectob{p} + \matmor{R}_{\monelb}\vectmor{t}_{\monela} + \vectmor{t}_{\monelb}.
    \end{aligned}
\end{equation}

Alternatively, we can first compose the transformation~$\monela$ with~$\monelb$:
%
\begin{equation}
    \label{eq:rototranslation-series-2}
    \tupp{\matmor{R}_{\monela}, \vectmor{t}_{\monela}} \mtimesof{\mgSEn} \tupp{\matmor{R}_{\monelb}, \vectmor{t}_{\monelb}} = \tupp{\matmor{R}_{\monelb} \matmor{R}_{\monela}, \matmor{R}_{\monelb} \vectmor{t}_{\monela} + \vectmor{t}_{\monelb}},
\end{equation}
%
and then apply this composite transformation to the point~$\vectob{p}$:
\begin{equation}
    \label{eq:rototranslation-series-3}
    \mgact(\tupp{\matmor{R}_{\monelb} \matmor{R}_{\monela}, \matmor{R}_{\monelb} \vectmor{t}_{\monela} + \vectmor{t}_{\monelb}}, \vectob{p})
    = \matmor{R}_{\monelb}\matmor{R}_{\monela}\vectob{p} + \matmor{R}_{\monelb}\vectmor{t}_{\monela} + \vectmor{t}_{\monelb}.
\end{equation}
%
Thus, we have proved this property
%
\begin{equation}
    \label{eq:rototranslation-action}
    \mgact(\monelb, \mgact(\monela, \vectob{p})) = \mgact(\monela \mtimes \monelb, \vectob{p}),
\end{equation}
which is graphically reported in \cref{fig:graphical-roto-action}.

\begin{figure}[h]
    \includesag{10_graph_roto_action}
    \caption{Graphical representation of roto-translation action.}
    \label{fig:graphical-roto-action}
\end{figure}
%
The notion of \SY{semigroup action} generalizes this property.
}

\todotextjira{26}{\bernina: @JL: add more exercises}

\todotext{graded exercise commented out here}
%\begin{gradedexercise}[\exname{UncurryingGroupActions}]
%    \label{ex:UncurryingGroupActions}
%    Prove \cref{lem:uncurrying-group-actions}.
%\end{gradedexercise}
%\solutionof{UncurryingGroupActions}

\begin{gradedexercise}[\exname{MatrixMultAction}]
    \label{ex:MatrixMultAction}
    Let~$\setA = \reals^n$, and let~$\mgGLn$ be the group of invertible~$\ntimesn$ matrices.
    Let
    \begin{equation}
        \label{eq:action-matrix-mult}
        \defmapset{
            \alpha
        }{
            \mgGLn \cartprod \setA
        }{
            \setA
        }{
            \tup{\matmor{M}, \vectmor{v}}
        }{
            \matmor{M}\vectmor{v}
        }
    \end{equation}
    be the usual multiplication of matrices with vectors.
    Check that \cref{eq:action-matrix-mult} defines a left group action operation of~$\mgGLn$ on~\setA.
\end{gradedexercise}

\solutionof{MatrixMultAction}



